\documentclass[11pt,a4paper]{article}
\usepackage[utf8]{inputenc}
\usepackage[T1]{fontenc}
\usepackage[margin=0.75in]{geometry}
\usepackage{enumitem}
\usepackage{titlesec}
\usepackage{xcolor}
\usepackage{hyperref}

\definecolor{primarycolor}{RGB}{33,37,41}
\definecolor{secondarycolor}{RGB}{0,102,204}
\definecolor{accentcolor}{RGB}{108,117,125}
\hypersetup{colorlinks=true,urlcolor=secondarycolor}
\titleformat{\section}{\Large\bfseries\color{primarycolor}}{}{0em}{}
\titlespacing*{\section}{0pt}{1.1em}{0.45em}
\setlist[itemize]{leftmargin=*,itemsep=0.2em,topsep=0.25em}
\pagestyle{empty}

\begin{document}
\begin{center}
{\Huge\bfseries\color{primarycolor} Tobias Mwangi-Hartley}\\[0.3em]
{\large\color{accentcolor} Data Scientist | Applied Machine Learning}\\[0.3em]
Edinburgh, UK | \href{mailto:tobias.mwangihartley@example.com}{tobias.mwangihartley@example.com} | +44 7700 900110 | \href{https://linkedin.com/in/tmh-example}{linkedin.com/in/tmh-example}
\end{center}

\section*{Professional Summary}
Data scientist with 6 years in analytics and data science, including 4+ years as a data scientist, focused on experimentation, forecasting and customer analytics in a subscription business. Built operational forecasting and churn-scoring workflows using Python, BigQuery and dbt, with measurable business outcomes. Strong statistical foundation and experience translating business questions into analyses people act on; now moving toward machine learning engineering.

\section*{Experience}
\textbf{Data Scientist | Thistlebox, meal-kit subscriptions | Edinburgh} \hfill \textit{Feb 2022–Present}
\begin{itemize}
\item Design and analyse approximately 40 A/B tests per year for growth and product teams; built a sequential testing calculator used by approximately 25 colleagues.
\item Built a six-week-ahead weekly recipe and region demand forecasting workflow using LightGBM and calendar, price and menu-position features. A scheduled notebook writes forecasts to BigQuery for operations to use in Looker; WAPE decreased from 23\% to 14\%, with estimated food waste reduction of 9\%.
\item Built an XGBoost churn propensity model scored weekly in BigQuery to target retention offers to top-decile customers. A 10\% holdout validated a 2.1 percentage-point churn reduction in the treated group.
\item Write production SQL in BigQuery and maintain dbt models for the analytics layer.
\item Since 2025, have written pytest unit tests for forecasting feature code, reaching approximately 60\% coverage.
\item Built a Streamlit app for planning and analysing sequential A/B tests using statsmodels; wrote its Dockerfile. The data platform team handles deployment to Cloud Run.
\item Present results to leadership monthly and have delivered internal “Stats for PMs” training three times.
\end{itemize}

\textbf{Junior Data Analyst | Firth Energy Retail | Glasgow} \hfill \textit{Sep 2020–Jan 2022}
\begin{itemize}
\item Conducted customer segmentation and tariff analysis using R and SQL.
\item Automated monthly regulator reports with R Markdown, saving the team approximately two days each month.
\end{itemize}

\section*{Additional Project}
\textbf{Image Classification Course Project} \hfill \textit{2025}
\begin{itemize}
\item Fine-tuned a small CNN in PyTorch on a public food-images dataset, achieving approximately 80\% accuracy; a personal learning project completed as part of an online deep learning course.
\end{itemize}

\section*{Education}
\textbf{MSc Statistics | University of the Lothians | 2020 | Distinction}\\
Dissertation: Bayesian hierarchical models for small-area energy demand\\[0.35em]
\textbf{BSc Mathematics | University of the Lothians | 2019}

\section*{Technical Skills}
Python (pandas, scikit-learn, LightGBM, XGBoost, statsmodels, Streamlit); R (tidyverse, brms); SQL (BigQuery); dbt; Git; pytest; Docker; Vertex AI Workbench scheduled notebooks; Looker; A/B testing; causal inference; feature engineering; time-series forecasting; Bayesian statistics; PyTorch (course project).

\section*{Languages}
English (native); Swahili (conversational)
\end{document}